\documentclass[a4paper,12pt]{article}
\usepackage[utf8]{inputenc}
\usepackage[T1]{fontenc}
\usepackage[margin=2.5cm]{geometry}
\usepackage{ctex} % 支持中文
\usepackage{setspace} % 设置行距
\usepackage{titlesec} % 标题格式
\usepackage{xcolor} % 颜色支持

% 设置行距，由于是纯文本墙，稍微增加行距以防完全无法阅读
\onehalfspacing

% 设置标题格式
\titleformat{\section}{\Large\bfseries\color{blue!40!black}}{}{0em}{}

% 元数据
\title{\textbf{2024年度环境、社会及管治（ESG）报告节选}}
\author{公司董事会办公室}
\date{2024年3月}

\begin{document}

\maketitle

\section{1. 前言：服务国家战略，守护自贸港门户}

2024年对于我们而言是极不平凡的一年，这一年是全省推进自贸港封关运作的攻坚之年，也是我们作为本土航空基础设施平台迎接新机遇的关键之年。本公司始终牢记“外防输入、内保畅通”的使命，紧紧围绕“一主一次，两支一货运”的机场群战略布局，全力提升枢纽能级并服务于国家开放战略。在这一年里，我们不仅关注飞机的起降与旅客的输送，更将目光投向了更长远的未来，思考如何让蓝天更蓝，如何让旅客的笑容更灿烂，以及如何让企业的根基更稳固。我们致力于通过绿色低碳运营、真情服务与稳健治理，为利益相关方创造可持续价值。

\section{2. 环境范畴：数据里的绿色承诺}

我们深知绿色是自贸港最鲜明的底色，因此在2024年积极响应国家“双碳”目标，通过技术创新与精细化管理显著降低了运营过程中的碳足迹。为了让公众清晰了解我们的环境影响，我们对全年的排放与能耗数据进行了详尽的统计与披露。在温室气体排放方面，2024年公司温室气体排放总量控制在51,518.71吨二氧化碳当量，其中范围1（直接排放，如汽油燃烧等）为11,214.39吨二氧化碳当量，范围2（间接排放，主要来自外购电力）为40,304.32吨二氧化碳当量，折算后的碳排放强度为0.12吨二氧化碳当量/万元。在能源消耗的具体细节上，全年综合能源消耗量折合标准煤13,538.22吨，其中外购电力消耗为7,511.05万千瓦时，汽油消耗为19.10万升。与此同时，我们非常重视水资源管理与废弃物处理，全年水资源消耗总量为145.94万吨，通过循环利用系统实现了中水回用量13.61万吨，让珍贵的水资源得到了二次生命，而在废弃物管理方面，产生的6,230.58吨无害废弃物和0.58吨有害废弃物全部实现了100\%的合规处理率。

除了基础数据的管控，我们还致力于打造“双碳机场”，其中的核心举措就是推广APU替代设施。传统的机场运行中，飞机停在地面时如果通过燃烧航空燃油来维持空调和照明（即使用APU辅助动力装置），会产生巨大的噪音和碳排放，为了解决这个问题，我们在旗下核心门户机场投入使用了31套APU替代设备（包括近机位18套和远机位13套），这就好比让飞机在地面“插电”休息，直接使用地面的电力设备供电供气而不再消耗燃油，这项技术全年累计减少碳排放25,382吨二氧化碳当量，旗下核心门户机场也因此荣获了“双碳机场”二星级称号。此外，我们充分利用海南得天独厚的光热资源，在旗下核心商业综合体的屋顶建设光伏项目，年均发电量达到450万千瓦时，占其总用电量的10\%，年减排2,263.59吨二氧化碳，旗下临空产业园的屋顶光伏装机容量更是达到了8.6MWp，预计年减排25.75万吨。在交通运输方面，核心机场新能源车占比已达43.6\%（共221辆），并新建充电桩约100个，年减碳761吨。在绿色建筑领域，在建的省会城市地标项目“海南中心”（双子塔）获得LEED CS金级预认证，某新区F07项目也获得了绿色建筑一星级标准。

\section{3. 社会范畴：从应急风暴到温情服务}

机场不仅仅是交通枢纽，更是连接人与人的温情纽带，2024年我们在保障安全、提升服务、关爱员工和回馈社会方面交出了务实的答卷。首先在运营绩效方面，随着旅游市场的强劲复苏，旗下控股及管理输出的9家机场总体起降达到16.11万架次，旅客吞吐量共计2,462.62万人次，其中核心门户机场旅客吞吐量重回“2000万级”俱乐部（位列全国第29位），重点区域机场也开通了至吉隆坡的首条国际客运航线。为了提升旅客体验，我们实施了“一机双屏”海关安检融合查验模式，让旅客只需一次排队即可完成两项检查，同时在核心机场启用了“一站式综合服务中心”，为外籍旅客提供支付绑定、通信文旅等服务，真正实现落地即融入。

在安全应急方面，2024年9月超强台风“摩羯”的正面袭击是对我们最大的考验。台风登陆前，我们迅速启动“防汛防风专项应急预案”，运管委提前进行航班动态调整，对机坪设施进行地毯式加固。台风过境后，全体员工展现了惊人的“海南速度”，迅速清理跑道异物并修复受损设施，在9月3日至7日期间，核心机场迅速恢复运行并累计执行航班1,300架次，有力保障了进出岛通道畅通。在日常安全管理上，我们上线了“安全智控”平台，关联制度手册9,626项，全年生成电子检查单31万份，确保安全标准落实到每一个动作，全年实现一般征候万架次率为0，责任事故为0，安全培训覆盖率100\%。

在员工发展与社会贡献方面，公司现有员工10,611人，其中女性占比38.9\%（4,133人），少数民族员工179人，董事会女性占比22\%。我们全年投入培训资金1,119.54万元，人均培训时长达到105.88小时。与此同时，我们积极回馈社会，全年对外捐赠资金15.48万元，投入乡村振兴资金190万元，志愿服务参与人次超过10,000人次，总时长达500,000小时，并圆满完成了博鳌亚洲论坛年会、消博会、少数民族运动会等重大保障任务。

\section{4. 管治范畴：合规体系与经济责任}

透明高效的治理是企业可持续发展的基石，我们构建了规范的“五会一层”治理体系，董事会由9人组成，其中独立董事3人，全年共召开股东大会5次、董事会13次，并披露公告70份，确保信息公开透明。在合规风控方面，我们构建了“三全三化”大监督体系，梳理风险点358个，反腐败培训覆盖了95\%的管理层和90\%的员工，培训总时长400小时，且供应商廉洁承诺书签署率达到100\%，其中包括对23家工程类供应商的严格遴选。作为自贸港建设的参与者，我们还积极通过参股布局离岛免税业务，2024年公司间接参与的5家离岛免税店线下销售额约52亿元，市场占比约17\%，为地方经济发展做出了实质性贡献。

\end{document}